\documentclass[bachelor, och, pract]{SCWorks}
% параметр - тип обучения - одно из значений:
%    spec     - специальность
%    bachelor - бакалавриат (по умолчанию)
%    master   - магистратура
% параметр - форма обучения - одно из значений:
%    och   - очное (по умолчанию)
%    zaoch - заочное
% параметр - тип работы - одно из значений:
%    referat    - реферат
%    coursework - курсовая работа (по умолчанию)
%    diploma    - дипломная работа
%    pract      - отчет по практике
%    pract      - отчет о научно-исследовательской работе
%    autoref    - автореферат выпускной работы
%    assignment - задание на выпускную квалификационную работу
%    review     - отзыв руководителя
%    critique   - рецензия на выпускную работу
% параметр - включение шрифта
%    times    - включение шрифта Times New Roman (если установлен)
%               по умолчанию выключен
\usepackage[T2A]{fontenc}
\usepackage[cp1251]{inputenc}
\usepackage{graphicx}
\usepackage[final]{pdfpages}

\usepackage[sort,compress]{cite}
\usepackage{amsmath}
\usepackage{amssymb}
\usepackage{amsthm}
\usepackage{fancyvrb}
\usepackage{longtable}
\usepackage{array}
\usepackage[english,russian]{babel}
\usepackage{caption}
\captionsetup[figure]{font= normalsize, labelfont=normalsize}


\usepackage[colorlinks=true]{hyperref}

\usepackage{caption}
\captionsetup[figure]{font= normalsize, labelfont=normalsize}


\begin{document}

% Кафедра (в родительном падеже)
\chair{дискретной математики и информационных технологий}

% Тема работы
\title{Основы работы с \LaTeX}

% Курс
\course{2}

% Группа
\group{221}

% Факультет (в родительном падеже) (по умолчанию "факультета КНиИТ")
%\department{факультета КНиИТ}

% Специальность/направление код - наименование
%\napravlenie{02.03.02 "--- Фундаментальная информатика и информационные технологии}
%\napravlenie{02.03.01 "--- Математическое обеспечение и администрирование информационных систем}
\napravlenie{09.03.01 "--- Информатика и вычислительная техника}
%\napravlenie{09.03.04 "--- Программная инженерия}
%\napravlenie{10.05.01 "--- Компьютерная безопасность}

% Для студентки. Для работы студента следующая команда не нужна.
%\studenttitle{студентки}

% Фамилия, имя, отчество в родительном падеже
\author{Мусатов Федор Алексеевич}

% Заведующий кафедрой
\chtitle{доцент, к.\,ф.-м.\,н.} % степень, звание
\chname{Л.\,Б.\,Тяпаев}

%Научный руководитель (для реферата преподаватель проверяющий работу)
\satitle{ст. преподаватель} %должность, степень, звание
\saname{М.\,В.\,Белоконь}


% Семестр (только для практики, для остальных
% типов работ не используется)
\term{3}

% Наименование практики (только для практики, для остальных
% типов работ не используется)
\practtype{учебная}

% Продолжительность практики (количество недель) (только для практики,
% для остальных типов работ не используется)
\duration{18}

% Даты начала и окончания практики (только для практики, для остальных
% типов работ не используется)
\practStart{01.09.2024}
\practFinish{31.12.2024}

% Год выполнения отчета
\date{2024}

\maketitle

% Включение нумерации рисунков, формул и таблиц по разделам
% (по умолчанию - нумерация сквозная)
% (допускается оба вида нумерации)
%\secNumbering


\tableofcontents


% Раздел "Обозначения и сокращения". Может отсутствовать в работе
%\abbreviations
%\begin{description}
%    \item $|A|$  "--- количество элементов в конечном множестве $A$;
%    \item $\det B$  "--- определитель матрицы $B$;
%    \item ИНС "--- Искусственная нейронная сеть;
%    \item FANN "--- Feedforward Artifitial Neural Network
%\end{description}

% Раздел "Определения". Может отсутствовать в работе
%\definitions

% Раздел "Определения, обозначения и сокращения". Может отсутствовать в работе.
% Если присутствует, то заменяет собой разделы "Обозначения и сокращения" и "Определения"
%\defabbr


% Раздел "Введение"
\intro
Целью данной практики является приобретение навыков оформления студенческой работы средствами системы компьютерной вёрстки \LaTeX \ [1].

Для выполнения поставленной цели необходимо решить следующие задачи:
\begin{itemize}
    \item ознакомиться со стандартом СТО 1.04.01 -- 2019 <<КУРСОВЫЕ РАБОТЫ (ПРОЕКТЫ) И ВЫПУСКНЫЕ КВАЛИФИКАЦИОННЫЕ РАБОТЫ. ПОРЯДОК ВЫПОЛНЕНИЯ, СТРУКТУРА И ПРАВИЛА ОФОРМЛЕНИЯ>> \ [2];
    \item изучить основы создания документов в \LaTeX: создание структуры документа, набор и форматирование текста, создание формул, вставку изображений и таблиц;
    \item освоить работу с шаблоном для оформления студенческих работ, предоставленным факультетом, сверстать с его помощью реферат и отчёт о практике.
\end{itemize}

Основы работы с \LaTeX изложены в источниках \ [3-5].

\section{Реферат по теме <<Тема реферата>>}

В данном разделе показан результат компиляции, исходный код документа представлен в приложении А.

\includepdf[pages=-]{Referat.pdf}

% Раздел "Заключение"
\conclusion
В процессе прохождения практики я ознакомился со стандартом оформления студенческой работы, мной были изучены основы работы с системой компьютерной вёрстки \LaTeX,  свёрстаны реферат и отчёт об учебной практике.

%Библиографический список, составленный вручную, без использования BibTeX
%
\begin{thebibliography}{99}

  \bibitem {1}LaTeX -- Википедия [Электронный ресурс] URL: https://ru.wikipedia.org/wiki/LaTeX (дата обращения: 23.12.2022
  \label{1}
  \bibitem {2}СТО 1.04.01 -- 2019 <<КУРСОВЫЕ РАБОТЫ (ПРОЕКТЫ) И ВЫПУСКНЫЕ КВАЛИФИКАЦИОННЫЕ РАБОТЫ. ПОРЯДОК ВЫПОЛНЕНИЯ, СТРУКТУРА И ПРАВИЛА ОФОРМЛЕНИЯ>> [Электронный ресурс] URL:https://www.sgu.ru/sites/default/files/documents/2021/oformlenie\_\\kursovyh\_i\_diplomnyh\_rabot.pdf (дата обращения: 23.12.2022)
  \label{2}
  \bibitem {3}Воронцов К.В. LaTex2e в примерах [Электронный ресурс] URL: http://www.ccas.ru/voron/download/voron05latex.pdf (дата обращения: 23.12.2022)
  \label{3}
  \bibitem {4}Столяров А.В. Сверстай диплом красиво: LaTeX за три дня Электронный ресурс] URL: http://www.stolyarov.info/books/pdf/latex3days.pdf (дата обращения: 23.12.2022)
  \label{4}
  \bibitem {5}LaTeX  -- Wikibooks, open books for an open world [Электронный ресурс] URL: https://en.wikibooks.org/wiki/LaTeX (дата обращения: 23.12.2022)
  \label{5}
 \bibitem {6}источник из реферата
  \label{6}
 \bibitem {7}источник из реферата
  \label{7}
 \bibitem {8}источник из реферата
  \label{8}
 \bibitem {9}источник из реферата
  \label{9}
 \bibitem {10}источник из реферата
   \label{10}
 
\end{thebibliography}

%Библиографический список, составленный с помощью BibTeX
%
%\bibliographystyle{gost780uv}
%\bibliography{thesis}

% Окончание основного документа и начало приложений
% Каждая последующая секция документа будет являться приложением
\appendix

\section{Исходный код документа}

\begin{Verbatim}[fontsize=\small, numbers=left]
\documentclass[bachelor, och, referat]{SCWorks}
% параметр - тип обучения - одно из значений:
%    spec     - специальность
%    bachelor - бакалавриат (по умолчанию)
%    master   - магистратура
% параметр - форма обучения - одно из значений:
%    och   - очное (по умолчанию)
%    zaoch - заочное
% параметр - тип работы - одно из значений:
%    referat    - реферат
%    coursework - курсовая работа (по умолчанию)
%    diploma    - дипломная работа
%    pract      - отчет по практике
%    pract      - отчет о научно-исследовательской работе
%    autoref    - автореферат выпускной работы
%    assignment - задание на выпускную квалификационную работу
%    review     - отзыв руководителя
%    critique   - рецензия на выпускную работу
% параметр - включение шрифта
%    times    - включение шрифта Times New Roman (если установлен)
%               по умолчанию выключен
\usepackage[T2A]{fontenc}
\usepackage[cp1251]{inputenc}
\usepackage{graphicx}

\usepackage[sort,compress]{cite}
\usepackage{amsmath}
\usepackage{amssymb}
\usepackage{amsthm}
\usepackage{fancyvrb}
\usepackage{longtable}
\usepackage{array}
\usepackage[english,russian]{babel}
\usepackage{caption}
\captionsetup[figure]{font= normalsize, labelfont=normalsize}


\usepackage[colorlinks=true]{hyperref}

\usepackage{caption}
\captionsetup[figure]{font= normalsize, labelfont=normalsize}


\begin{document}
	
	% Кафедра (в родительном падеже)
	\chair{дискретной математики и информационных технологий}
	
	% Тема работы
	\title{Тема реферата}
	
	% Курс
	\course{2}
	
	% Группа
	\group{221}
	
	% Факультет (в родительном падеже) (по умолчанию "факультета КНиИТ")
	%\department{факультета КНиИТ}
	
	% Специальность/направление код - наименование
	%\napravlenie{02.03.02 "--- Фундаментальная информатика и информационные технологии}
	%\napravlenie{02.03.01 "--- Математическое обеспечение и администрирование информационных систем}
	\napravlenie{09.03.01 "--- Информатика и вычислительная техника}
	%\napravlenie{09.03.04 "--- Программная инженерия}
	%\napravlenie{10.05.01 "--- Компьютерная безопасность}
	
	% Для студентки. Для работы студента следующая команда не нужна.
	%\studenttitle{Студентки}
	
	% Фамилия, имя, отчество в родительном падеже
	\author{Иванова Ивана Ивановича}
	
	% Заведующий кафедрой
	\chtitle{доцент, к.\,ф.-м.\,н.} % степень, звание
	\chname{Л.\,Б.\,Тяпаев}
	
	%Научный руководитель (для реферата преподаватель проверяющий работу)
	\satitle{старший преподаватель} %должность, степень, звание
	\saname{М.\,В.\,Белоконь}
	
	
	% Семестр (только для практики, для остальных
	% типов работ не используется)
	\term{3}
	
	% Наименование практики (только для практики, для остальных
	% типов работ не используется)
	\practtype{Преддипломная}
	
	% Продолжительность практики (количество недель) (только для практики,
	% для остальных типов работ не используется)
	\duration{4}
	
	% Даты начала и окончания практики (только для практики, для остальных
	% типов работ не используется)
	\practStart{30.04.2021}
	\practFinish{27.05.2021}
	
	% Год выполнения отчета
	\date{2023}
	
	\maketitle
	
	% Включение нумерации рисунков, формул и таблиц по разделам
	% (по умолчанию - нумерация сквозная)
	% (допускается оба вида нумерации)
	%\secNumbering
	
	
	\tableofcontents
	
	
	% Раздел "Обозначения и сокращения". Может отсутствовать в работе
	\abbreviations
	\begin{description}
		\item $|A|$  "--- количество элементов в конечном множестве $A$;
		\item $\det B$  "--- определитель матрицы $B$;
		\item ИНС "--- Искусственная нейронная сеть;
		\item FANN "--- Feedforward Artifitial Neural Network
	\end{description}
	
	% Раздел "Определения". Может отсутствовать в работе
	%\definitions
	
	% Раздел "Определения, обозначения и сокращения". Может отсутствовать в работе.
	% Если присутствует, то заменяет собой разделы "Обозначения и сокращения" и "Определения"
	%\defabbr
	
	
	% Раздел "Введение"
	\intro
	Целью настоящей работы является создание примера оформления студенческой работы средствами системы \LaTeX.
	
	Поставлена задача оформить документ в соответствии:
	\begin{itemize}
		\item со стандартом СТО 1.04.01-2012 Порядком выполнения, структурой и правилами оформления курсовых работ (проектов)
		и выпускных квалификационных работ, принятых в Саратовском государственном университете в 2012 году;
		\item с правилами оформления титульного листа отчета о прохождении практики в соответствии со стандартом СТО 1.01-2005.
	\end{itemize}
	
	Изложенный ниже текст не имеет особого смысла и приведен только для демонстрации оформления своих элементов.
	
	\section{Пример оформления текста}
	\subsection{Пример основных элементов математического текста}
	
	Внутритекстовая формула $\frac{1}{\varepsilon^*}=\frac{1}{\varepsilon_\infty}-\frac{1}{\varepsilon_0}$.
	Пример одиночной ссылки на литературу~\cite{1}. Пример множественной ссылки на
	литературу~\cite{2}
	.
	\begin{equation}
		F(x)=\int\limits_a^bf(x)\,dx.
	\end{equation}
	
	Ссылка на рисунок~\ref{fig:f}.
	\begin{figure}[!ht]
		\centering
		\includegraphics[width=6cm]{101.png}
		\caption{\label{fig:f}%
			Подпись к рисунку}
	\end{figure}
	
	
	
	Если разность энергий электронно"=дырочных уровней $E_2-E_1$ близка к энергии предельного оптического фонона $\hbar\Omega_{LO}$, то в разложении волновых функций полного гамильтониана можно ограничиться нулевым приближением для всех состояний, за исключением близких по значению к $E_2$. Волновые функции последних представляют собой следующие комбинации почти вырожденных состояний\cite{3}.
	
	\subsection{Еще элементы математического текста}
	Нейрон является составной частью нейронной сети. Он состоит из
	элементов трех типов: умножителей (синапсов), сумматора и
	нелинейного преобразователя. Синапсы осуществляют связь между
	нейронами, умножают входной сигнал на число, характеризующее силу
	связи (вес синапса). Сумматор выполняет сложение сигналов,
	поступающих по синаптическим связям от других нейронов, и внешних
	входных сигналов. Нелинейный преобразователь реализует нелинейную
	функцию одного аргумента "--- выхода сумматора. Эта функция
	называется функцией активации или передаточной функцией. На рисунке~\ref{neuron} приведено строение одного нейрона.
	
	Нейрон в целом реализует скалярную функцию векторного аргумента.
	Математическая модель нейрона:
	\[
	s = \sum\limits_{i = 1}^n w_i x_i + b,
	\]
	\[
	y = f(s),
	\]
	где $w_i $ "--- вес синапса; $i = 1,\ldots ,n$; $b$ "--- значение
	смещения; $s$ "--- результат суммирования; $x_i $ "--- $i$-тый
	компонент входного вектора (входной сигнал), \linebreak $i = 1,\ldots, n$;
	$y$ "--- выходной сигнал нейрона; $n$ "--- число входов нейрона;
	$f(s)$ "--- нелинейное преобразование (функция активации).
	\begin{figure}[ht]
		\centering
		\includegraphics{101.png}
		\caption{Нейрон}\label{neuron}
	\end{figure}
	
	В качестве функции активации нейронов берут обычно одну из
	следующих:
	\begin{itemize}
		\item пороговая функция активации;
		\item экспоненциальная сигмоида;
		\item рациональная сигмоида;
		\item гиперболический тангенс.
	\end{itemize}
	
	Данные функции активации обладают таким важным свойством как
	нелинейность. Нелинейность функции активации принципиальна для
	построения нейронных сетей. Если бы нейроны были линейными
	элементами, то любая последовательность нейронов также производила
	бы линейное преобразование и вся нейронная сеть была бы
	эквивалентна одному нейрону (или одному слою нейронов в случае
	нескольких выходов). Нелинейность разрушает суперпозицию и
	приводит к тому, что возможности нейросети существенно выше
	возможностей отдельных нейронов.
	
	\subsection{Снова математический текст}
	Опишем самую популярную архитектуру
	"--- многослойный персептрон с последовательными связями и
	сигмоидальной функцией активации (\foreignlanguage{english}{Feedforward Artifitial Neural
		Network, FANN}).
	
	В многослойных нейронных сетях с последовательными связями нейроны
	делятся на группы с общим входным сигналом "--- слои. Стандартная
	сеть состоит из $L$ слоев, пронумерованных слева направо. Каждый
	слой содержит совокупность нейронов с едиными входными сигналами.
	Внешние входные сигналы подаются на входы нейронов входного слоя
	(его часто нумеруют как нулевой), а выходами сети являются
	выходные сигналы последнего слоя. Кроме входного и выходного слоев
	в многослойной нейронной сети есть один или несколько скрытых
	слоев, соединенных последовательно в прямом направлении и не
	содержащих связей между элементами внутри слоя и обратных связей
	между слоями. Число нейронов в слое может быть любым и не зависит
	от количества нейронов в других слоях. Архитектура нейронной сети
	прямого распространения сигнала приведена на рисунке~\ref{net1}.
	
	На каждый нейрон первого слоя подаются все элементы внешнего
	входного сигнала. Все выходы нейронов $i$-го слоя подаются на
	каждый нейрон слоя $i+1$.
	
	Нейроны выполняют взвешенное суммирование элементов входных
	сигналов. К сумме прибавляется смещение нейрона. Над результатом
	суммирования выполняется нелинейное преобразование "--- функция
	активации (передаточная функция). Значение функции активации есть
	выход нейрона. Приведем схему многослойного персептрона. Нейроны
	представлены кружками, связи между нейронами "--- линиями со
	стрелками.
	
	\begin{figure}[ht]
		\centering
		\includegraphics{101.png}
		\caption{Архитектура многослойной сети прямого
			распространения}\label{net1}
	\end{figure}
	
	Функционирование сети выполняется в соответствии с формулами:
	\[
	s_j^{\left[ k \right]} = \sum\limits_{i = 1}^{N_{k - 1} }
	{w_{ji}^{\left[ k \right]} y_i^{\left[ {k - 1} \right]} +
		b_j^{[k]} ,\ \ j = 1,\ldots ,N_k ,\ \ k = 1,\ldots ,L;}
	\]
	\[
	y_j^{\left[ k \right]} = f(s_j^{\left[ k \right]} ),\ \ j =
	1,\ldots ,N_k ,\ \ k = 1,\ldots ,L-1,
	\]
	\[
	y_j^{\left[ L \right]} = s_j^{\left[ L \right]} ,
	\]
	где
	\begin{itemize}
		\item
		$y_i^{\left[ {k - 1} \right]}$ "--- выходной сигнал $i$-го нейрона
		$(k - 1)$-го слоя; \item $w_{ji}^{\left[ k \right]}$ "--- вес связи
		между  $j$-м нейроном слоя $(k-1)$  и  $i$-м нейроном $k$-го
		слоя;
		\item
		$b_j^{\left[ k \right]}$ "--- значение смещения $j$-го
		нейрона $k$-го слоя;
		\item
		$y = f(s)$ "--- функция активации;
		\item
		$y_j^{\left[ k \right]}$ "--- выходной сигнал $j$-го
		нейрона $k$-го слоя;
		\item
		$N_k$ "--- число узлов слоя $k$;
		\item
		$L$ "--- общее число основных слоев;
		\item
		$n = N_0$ "--- размерность входного вектора;
		\item
		$m = N_L$ "---
		размерность выходного вектора сети.
	\end{itemize}
	
	На рисунке~\ref{net2} представлена сеть прямого распространения
	сигнала с 5 входами, 3 нейронами в скрытом слое и 2 нейронами в
	выходном слое.
	\begin{figure}[hb]
		\centering
		\includegraphics{102.png}
		\caption{Пример нейронной сети}\label{net2}
	\end{figure}
	
	
	\section{Раздел с подразделами}
	\subsection{Текст с формулами и леммой}
	
	Обозначим $[y_0,y_1,\ldots,y_p;f]$ разделенную разность порядка $p$ функции $f$ по узлам $y_0<y_1<\ldots<y_p$.
	
	Обозначим $L_pf(x;y_0,y_1,\ldots,y_p)$ интерполяционный полином Ньютона функции $f$ по узлам $y_0,y_1,\ldots,y_p$:
	
	
	\begin{proof}
		Возьмем $x \in [x_{p-(2k+1)},x_{p-2k}]$,\ \ $k=0,\ldots,\left[p/2
		\right]$.
		
		Из условия леммы следует, что
		\begin{displaymath}
			[x_0,\ldots,x_{p-(2k+1)},x,x_{p-2k},\ldots,x_p;f] \geqslant
			0,
		\end{displaymath}
		
		
		Из равенства
		\begin{equation*}
			\Delta_p f(x;x_0,\ldots,x_p)=(L_pf(x;x_0,\ldots,x_p)-f(x))
			\prod_{0\leqslant i<j\leqslant p}(x_j-x_i).
		\end{equation*}
		и следует, что
		\begin{displaymath}
			L_pf(x;x_0,\ldots,x_p)\geqslant f(x).
		\end{displaymath}
		
		С учетом условия леммы мы получаем утверждение.
	\end{proof}
	
	\subsection{Название другого подраздела}
	\subsubsection{Более мелкий подраздел}
	Если разность энергий электронно"=дырочных уровней $E_2-E_1$ близка к энергии предельного оптического фонона $\hbar\Omega_{LO}$, то в разложении волновых функций полного гамильтониана можно ограничиться нулевым приближением для всех состояний, за исключением близких по значению к $E_2$.
	
	\subsubsection{Текст с таблицей}
	В таблице~\ref{table-1} представлены результаты сокращения словарей неисправностей для схем из каталога ISCAS'89.
	
	\begin{table}[!ht]
		\small
		\caption{Результат сокращения словарей неисправностей при помощи масок} \label{table-1}
		\begin{tabular}{|l|c|c|c|c|r|r|r|}
			\hline 1 & 2& 3& 4& 5& 6& 7& 8\\
			\hline S298 & 177 & 1932 & 341964 & 61 & 10797 & 3,16\% & 0,61\\
			\hline S344 & 240 & 1397 & 335280 & 59 & 14160 & 4,22\% & 0,53\\
			\hline S349 & 243 & 1474 & 358182 & 62 & 15066 & 4,21\% & 0,60\\
			\hline S382 & 190 & 12444 & 2364360 & 55 & 10450 & 0,44\% & 3,78\\
			\hline S386 & 274 & 2002 & 548548 & 91 & 24934 & 4,55\% & 1,40\\
			\hline S400 & 194 & 13284 & 2577096 & 58 & 11252 & 0,44\% & 4,28\\
			\hline S444 & 191 & 13440 & 2567040 & 60 & 11460 & 0,45\% & 4,26\\
			\hline S510 & 446 & 700 & 312200 & 70 & 31220 & 10,00\% & 0,63\\
			\hline S526 & 138 & 13548 & 1869624 & 38 & 5244 & 0,28\% & 2,41\\
			\hline S641 & 345 & 5016 & 1730520 & 132 & 45540 & 2,63\% & 7,06\\
			\hline S713 & 343 & 3979 & 1364797 & 131 & 44933 & 3,29\% & 5,61\\
			\hline S820 & 712 & 21185 & 15083720 & 244 & 173728 & 1,15\% & 126,99\\
			\hline S832 & 719 & 21603 & 15532557 & 253 & 181907 & 1,17\% & 135,18\\
			\hline S953 & 326 & 322 & 104972 & 91 & 29666 & 28,26\% & 0,27\\
			\hline S1423 & 293 & 750 & 219750 & 93 & 27249 & 12,40\% & 0,57\\
			\hline S1488 & 1359 & 22230 & 30210570 & 384 & 521856 & 1,73\% & 541,69\\
			\hline
		\end{tabular}
	\end{table}
	
	\subsubsection{Текст с кодом программы}
	Термин <<разреженная матрица>> впервые был предложен Гарри Марковицем. В 1989 он был награжден премией имени Джона фон Неймана в том числе и за вклад в теорию методов для разреженных матриц.
	
	В большинстве источников, разреженной матрицей называется матрица, в которой мало ненулевых элементов. Это нельзя назвать определением из-за слова <<мало>>. В понятие разреженной матрицы определяется так: <<Мы можем называть матрицу разреженной, если применение к ней методов, описываемых в книге, экономит память и/или время>>. Таким образом, следует дать определение алгоритму для разреженных матриц. Алгоритмом для разреженных матриц будем называть алгоритм, у которого время работы и необходимый объем памяти зависят от количества ненулевых элементов в матрице.
	
	Размерность квадратной матрицы $A$ будем обозначать $n$, а количество ненулевых элементов в ней $|A|$.
	
	Плотные матрицы обычно хранятся в качестве двумерного массива $n\times n$. Будем обозначать такой массив a. Разреженные матрицы не стоит хранить таким способом из-за слишком большого потребления памяти, которая будет занята в основном нулевыми элементами.
	
	Один из вариантов представления разреженных матриц в памяти компьютера "--- в виде трех массивов: \verb"column", \verb"value" и \verb"rowIndex". Размеры массивов \verb"column" и \verb"value" равны $|A|$. Размер \verb"rowIndex" равен $n+1$. Ненулевые элементы матрицы $A$ хранятся последовательно по строкам в этих массивах. Элемент \verb"column[i]" содержит номер столбца, в котором содержится \verb"i"-й ненулевой элемент, а \verb"value[i]" "--- его величину. Массив \verb"rowIndex[i]" содержит в себе индекс первого ненулевого элемента \verb"i"-й строки. Все ненулевые элементы \verb"i"-й строки содержатся в массивах \verb"column" и \verb"value" в элементах с индексами от \verb"rowIndex[i]" по \verb"rowIndex[i + 1]-1". Для удобства полагают \verb"rowIndex["$n$\verb"]"$=|A|$.
	
	Для примера рассмотрим следующую матрицу:
	\[
	\left(
	\begin{matrix}
		1 & 0 & 5 & 0 & 0 \\
		0 & 2 & 7 & 4 & 0 \\
		0 & 0 & 1 & 0 & 0 \\
		9 & 6 & 0 & 3 & 0 \\
		0 & 0 & 3 & 0 & 5
	\end{matrix}
	\right)
	\]
	
	Массивы \verb"column", \verb"value" и \verb"rowIndex" для этой матрицы представлены в таблице~\ref{tab:mat-arrays}.
	\begin{table}[ht]\small
		\caption{Массивы \texttt{column}, \texttt{value} и \texttt{rowIndex}}\label{tab:mat-arrays}
		\begin{tabular}{|l|c|c|c|c|c|c|c|c|c|c|c|c|} \cline{2-13}
			\multicolumn{1}{c|}{} & 0 & 1 & 2 & 3 & 4 & 5 & 6 & 7 & 8 & 9 & 10 & 11 \\ \cline{2-13}\hline
			\verb"column" & 0 & 2 & 1 & 2 & 3 & 2 & 0 & 1 & 3 & 2 & 4 &  \\ \hline\hline
			\verb"value" & 1 & 5 & 2 & 7 & 4 & 1 & 9 & 6 & 3 & 3 & 5 &  \\ \hline\hline
			\verb"rowIndex" & 0 & 2 & 5 & 6 & 9 & 11 &  &  &  &  &  &  \\ \hline
		\end{tabular}
	\end{table}
	
	Неизвестный вектор и вектор правой части хранятся в виде массивов размера $n$. Массив неизвестного вектора обозначают \verb"x", а массив правой части "--- \verb"rhs".
	
	Рассмотрим пример алгоритма для разреженных матриц. Алгоритм решения СЛАУ, представленной нижнетреугольной матрицей \verb"a", можно реализовать двумя вложенными циклами по \verb"n":
	\begin{Verbatim}[fontsize=\small, numbers=left]
		for(int i = 0; i $<$ n; ++i){
			x[i] = rhs[i];
			for(int j = 0; j $<$ i; ++j)
			x[i] -= a[i][j] * x[j];
			x[i] /= a[i][i];
		}

	
	Но, если матрица \verb"a" хранится в разреженном виде, то в данном алгоритме можно проходить только по ненулевым элементам \verb"a":
	\begin{Verbatim}[fontsize=\small, numbers=left]
		for(int i = 0; i $<$ n; ++i){
			x[i] = rhs[i];
			for(int j = rowIndex[i]; j $<$ rowIndex[i + 1] - 1; ++j)
			x[i] -= value[j] * x[column[j]];
			x[i] /= value[rowIndex[i + 1] - 1];
		}

	В первом случае оценка времени работы будет $O(n^{2})$, а во втором $O(|A|)$.
	
	Методы для разреженных матриц основаны на следующих главных принципах:
	
	\begin{enumerate}
		\item  Хранятся только ненулевые элементы матрицы.
		\item  Выполняются только те преобразования, которые действительно что"=то изменяют. В примере не имеет смысла вычитать из \verb"x[i]" значение \verb"x[j]*a[i][j]", если \verb"a[i][j]" равно нулю.
		\item  Число <<новых элементов>>, возникающих, например, во время исключения Гаусса, стараются уменьшить путем перестановок строк и столбцов матрицы.
	\end{enumerate}
	
	Пример оформления рисунка с 3-мя секциями показан на рисунке ~\ref{5} (а, б, в).
	
	\begin{figure}[hb]
		\centering
		\includegraphics{103.png}
		\caption{Архитектура многослойной сети прямого
			распространения: a) название подрисунка а б) название подрисунка б в) название подрисунка в}\label{5}
	\end{figure}
	
	Пример оформления рисунка с 3-мя секциями показан на рисунке ~\ref{ris} (а, б, в).
	
	\begin{figure}[hb]
		\begin{minipage}[h]{0.32\linewidth}
			\center{\includegraphics[width=0.7\linewidth]{101.png}}
		\end{minipage}
		\hfill
		\begin{minipage}[h]{0.32\linewidth}
			\center{\includegraphics[width=0.7\linewidth]{101.png}}
		\end{minipage}
		\hfill
		\begin{minipage}[h]{0.32\linewidth}
			\center{\includegraphics[width=0.7\linewidth]{101.png}}
		\end{minipage}
		\begin{minipage}[h]{1\linewidth}
			\begin{tabular}{p{0.32\linewidth}p{0.32\linewidth}p{0.32\linewidth}}
				\centering а) & \centering б) & \centering в) \\
			\end{tabular}
		\end{minipage}
		\vspace*{-1cm}
		
		\caption{Пример оформления: a) подрисунка а, б)
			подрисунка б, в) подрисунка в.}
		\label{ris}
	\end{figure}
	
	% Раздел "Заключение"
	\conclusion
	В настоящей работы приведен пример оформления студенческой работы средствами системы \LaTeX.
	
	Показано, как можно оформить документ в соответствии:
	\begin{itemize}
		\item с правилами оформления курсовых и выпускных квалификационных работ, принятых в Саратовском государственном университете в 2012 году;
		\item с правилами оформления титульного листа отчета о прохождении практики в соответствии со стандартом.
	\end{itemize}
	%Библиографический список, составленный вручную, без использования BibTeX
	%
	\begin{thebibliography}{99}
		
		\bibitem {1} Машинное зрение. Что это и как им пользоваться? Обработка изображений оптического источника: [Электронный ресурс] URL: https://habr.com/ru/post/350918/(дата обращения: 20.04.2021)
		\label{1}
		\bibitem {2} Гудков, В. А. Исследование молекулярной и надмолекулярной структуры ряда жидкокристаллических полимеров / В. А. Гудков // Журн. структур, химии. 1991. Т. 32, № 4. С. 86-91. 
		\label{2}
		\bibitem{3} Обработка и анализ цифровых изображений с примерами на LabVIEW IMAQ Vision/ Визильтер Ю. В., Желтов С. Ю., Князь В. А., Ходарев А. Н., Моржин А. В. – М.: ДМК Пресс, 2009. 464 с.
		\label{3}
		
	\end{thebibliography}
	
	%Библиографический список, составленный с помощью BibTeX
	%
	%\bibliographystyle{gost780uv}
	%\bibliography{thesis}
	
	% Окончание основного документа и начало приложений
	% Каждая последующая секция документа будет являться приложением
	\appendix
	
	\section{Нумеруемые объекты в приложении}
	
	\newpage 
	
	\begin{figure}
		
		Реферат "Тема реферата" \ выполнен мною самостоятельно, и на все источники, имеющиеся в реферате, даны соответствующие ссылки.
		
	\end{figure}
	
	\begin{figure}
		\begin{flushright}
			\underline{\hspace{4cm}}  \quad  \underline{\hspace{4cm}} 
			
			\scriptsize подпись, дата     \quad     \qquad  \qquad  \qquad        \scriptsize инициалы, фамилия  \qquad  \qquad  \qquad 
		\end{flushright}
	\end{figure}
\end{document}

\end{Verbatim}

\newpage 

\begin{figure}

Отчет о практике выполнен мною самостоятельно, и на все источники, имеющиеся в отчете, даны соответствующие ссылки.

\end{figure}

\begin{figure}
\begin{flushright}
\underline{\hspace{4cm}}  \quad  \underline{\hspace{4cm}} 

 \scriptsize подпись, дата     \quad     \qquad  \qquad  \qquad        \scriptsize инициалы, фамилия  \qquad  \qquad  \qquad 
\end{flushright}
\end{figure}
\end{document}
